\IfFileExists{stacks-project.cls}{%
\documentclass{stacks-project}
}{%
\documentclass{amsart}
}

% For dealing with references we use the comment environment
\usepackage{verbatim}
\newenvironment{reference}{\comment}{\endcomment}
%\newenvironment{reference}{}{}
\newenvironment{slogan}{\comment}{\endcomment}
\newenvironment{history}{\comment}{\endcomment}

% For commutative diagrams we use Xy-pic
\usepackage[all]{xy}

% We use 2cell for 2-commutative diagrams.
\xyoption{2cell}
\UseAllTwocells

% We use multicol for the list of chapters between chapters
\usepackage{multicol}

% This is generally recommended for better output
\usepackage{lmodern}
\usepackage[T1]{fontenc}

% For cross-file-references
\usepackage{xr-hyper}

% Package for hypertext links:
\usepackage{hyperref}

% For any local file, say "hello.tex" you want to link to please
% use \externaldocument[hello-]{hello}
\externaldocument[introduction-]{introduction}
\externaldocument[conventions-]{conventions}
\externaldocument[sets-]{sets}
\externaldocument[categories-]{categories}
\externaldocument[topology-]{topology}
\externaldocument[sheaves-]{sheaves}
\externaldocument[sites-]{sites}
\externaldocument[stacks-]{stacks}
\externaldocument[fields-]{fields}
\externaldocument[algebra-]{algebra}
\externaldocument[brauer-]{brauer}
\externaldocument[homology-]{homology}
\externaldocument[derived-]{derived}
\externaldocument[simplicial-]{simplicial}
\externaldocument[more-algebra-]{more-algebra}
\externaldocument[smoothing-]{smoothing}
\externaldocument[modules-]{modules}
\externaldocument[sites-modules-]{sites-modules}
\externaldocument[injectives-]{injectives}
\externaldocument[cohomology-]{cohomology}
\externaldocument[sites-cohomology-]{sites-cohomology}
\externaldocument[dga-]{dga}
\externaldocument[dpa-]{dpa}
\externaldocument[sdga-]{sdga}
\externaldocument[hypercovering-]{hypercovering}
\externaldocument[schemes-]{schemes}
\externaldocument[constructions-]{constructions}
\externaldocument[properties-]{properties}
\externaldocument[morphisms-]{morphisms}
\externaldocument[coherent-]{coherent}
\externaldocument[divisors-]{divisors}
\externaldocument[limits-]{limits}
\externaldocument[varieties-]{varieties}
\externaldocument[topologies-]{topologies}
\externaldocument[descent-]{descent}
\externaldocument[perfect-]{perfect}
\externaldocument[more-morphisms-]{more-morphisms}
\externaldocument[flat-]{flat}
\externaldocument[groupoids-]{groupoids}
\externaldocument[more-groupoids-]{more-groupoids}
\externaldocument[etale-]{etale}
\externaldocument[chow-]{chow}
\externaldocument[intersection-]{intersection}
\externaldocument[pic-]{pic}
\externaldocument[weil-]{weil}
\externaldocument[adequate-]{adequate}
\externaldocument[dualizing-]{dualizing}
\externaldocument[duality-]{duality}
\externaldocument[discriminant-]{discriminant}
\externaldocument[derham-]{derham}
\externaldocument[local-cohomology-]{local-cohomology}
\externaldocument[algebraization-]{algebraization}
\externaldocument[curves-]{curves}
\externaldocument[resolve-]{resolve}
\externaldocument[models-]{models}
\externaldocument[functors-]{functors}
\externaldocument[equiv-]{equiv}
\externaldocument[pione-]{pione}
\externaldocument[etale-cohomology-]{etale-cohomology}
\externaldocument[proetale-]{proetale}
\externaldocument[relative-cycles-]{relative-cycles}
\externaldocument[more-etale-]{more-etale}
\externaldocument[trace-]{trace}
\externaldocument[crystalline-]{crystalline}
\externaldocument[spaces-]{spaces}
\externaldocument[spaces-properties-]{spaces-properties}
\externaldocument[spaces-morphisms-]{spaces-morphisms}
\externaldocument[decent-spaces-]{decent-spaces}
\externaldocument[spaces-cohomology-]{spaces-cohomology}
\externaldocument[spaces-limits-]{spaces-limits}
\externaldocument[spaces-divisors-]{spaces-divisors}
\externaldocument[spaces-over-fields-]{spaces-over-fields}
\externaldocument[spaces-topologies-]{spaces-topologies}
\externaldocument[spaces-descent-]{spaces-descent}
\externaldocument[spaces-perfect-]{spaces-perfect}
\externaldocument[spaces-more-morphisms-]{spaces-more-morphisms}
\externaldocument[spaces-flat-]{spaces-flat}
\externaldocument[spaces-groupoids-]{spaces-groupoids}
\externaldocument[spaces-more-groupoids-]{spaces-more-groupoids}
\externaldocument[bootstrap-]{bootstrap}
\externaldocument[spaces-pushouts-]{spaces-pushouts}
\externaldocument[spaces-chow-]{spaces-chow}
\externaldocument[groupoids-quotients-]{groupoids-quotients}
\externaldocument[spaces-more-cohomology-]{spaces-more-cohomology}
\externaldocument[spaces-simplicial-]{spaces-simplicial}
\externaldocument[spaces-duality-]{spaces-duality}
\externaldocument[formal-spaces-]{formal-spaces}
\externaldocument[restricted-]{restricted}
\externaldocument[spaces-resolve-]{spaces-resolve}
\externaldocument[formal-defos-]{formal-defos}
\externaldocument[defos-]{defos}
\externaldocument[cotangent-]{cotangent}
\externaldocument[examples-defos-]{examples-defos}
\externaldocument[algebraic-]{algebraic}
\externaldocument[examples-stacks-]{examples-stacks}
\externaldocument[stacks-sheaves-]{stacks-sheaves}
\externaldocument[criteria-]{criteria}
\externaldocument[artin-]{artin}
\externaldocument[quot-]{quot}
\externaldocument[stacks-properties-]{stacks-properties}
\externaldocument[stacks-morphisms-]{stacks-morphisms}
\externaldocument[stacks-limits-]{stacks-limits}
\externaldocument[stacks-cohomology-]{stacks-cohomology}
\externaldocument[stacks-perfect-]{stacks-perfect}
\externaldocument[stacks-introduction-]{stacks-introduction}
\externaldocument[stacks-more-morphisms-]{stacks-more-morphisms}
\externaldocument[stacks-geometry-]{stacks-geometry}
\externaldocument[moduli-]{moduli}
\externaldocument[moduli-curves-]{moduli-curves}
\externaldocument[examples-]{examples}
\externaldocument[exercises-]{exercises}
\externaldocument[guide-]{guide}
\externaldocument[desirables-]{desirables}
\externaldocument[coding-]{coding}
\externaldocument[obsolete-]{obsolete}
\externaldocument[fdl-]{fdl}
\externaldocument[index-]{index}

% Theorem environments.
%
\theoremstyle{plain}
\newtheorem{theorem}[subsection]{Theorem}
\newtheorem{proposition}[subsection]{Proposition}
\newtheorem{lemma}[subsection]{Lemma}

\theoremstyle{definition}
\newtheorem{definition}[subsection]{Definition}
\newtheorem{example}[subsection]{Example}
\newtheorem{exercise}[subsection]{Exercise}
\newtheorem{situation}[subsection]{Situation}

\theoremstyle{remark}
\newtheorem{remark}[subsection]{Remark}
\newtheorem{remarks}[subsection]{Remarks}

\numberwithin{equation}{subsection}

% Macros
%
\def\lim{\mathop{\mathrm{lim}}\nolimits}
\def\colim{\mathop{\mathrm{colim}}\nolimits}
\def\Spec{\mathop{\mathrm{Spec}}}
\def\Hom{\mathop{\mathrm{Hom}}\nolimits}
\def\Ext{\mathop{\mathrm{Ext}}\nolimits}
\def\SheafHom{\mathop{\mathcal{H}\!\mathit{om}}\nolimits}
\def\SheafExt{\mathop{\mathcal{E}\!\mathit{xt}}\nolimits}
\def\Sch{\mathit{Sch}}
\def\Mor{\mathop{\mathrm{Mor}}\nolimits}
\def\Ob{\mathop{\mathrm{Ob}}\nolimits}
\def\Sh{\mathop{\mathit{Sh}}\nolimits}
\def\NL{\mathop{N\!L}\nolimits}
\def\CH{\mathop{\mathrm{CH}}\nolimits}
\def\proetale{{pro\text{-}\acute{e}tale}}
\def\etale{{\acute{e}tale}}
\def\QCoh{\mathit{QCoh}}
\def\Ker{\mathop{\mathrm{Ker}}}
\def\Im{\mathop{\mathrm{Im}}}
\def\Coker{\mathop{\mathrm{Coker}}}
\def\Coim{\mathop{\mathrm{Coim}}}

% Boxtimes
%
\DeclareMathSymbol{\boxtimes}{\mathbin}{AMSa}{"02}

%
% Macros for moduli stacks/spaces
%
\def\QCohstack{\mathcal{QC}\!\mathit{oh}}
\def\Cohstack{\mathcal{C}\!\mathit{oh}}
\def\Spacesstack{\mathcal{S}\!\mathit{paces}}
\def\Quotfunctor{\mathrm{Quot}}
\def\Hilbfunctor{\mathrm{Hilb}}
\def\Curvesstack{\mathcal{C}\!\mathit{urves}}
\def\Polarizedstack{\mathcal{P}\!\mathit{olarized}}
\def\Complexesstack{\mathcal{C}\!\mathit{omplexes}}
% \Pic is the operator that assigns to X its picard group, usage \Pic(X)
% \Picardstack_{X/B} denotes the Picard stack of X over B
% \Picardfunctor_{X/B} denotes the Picard functor of X over B
\def\Pic{\mathop{\mathrm{Pic}}\nolimits}
\def\Picardstack{\mathcal{P}\!\mathit{ic}}
\def\Picardfunctor{\mathrm{Pic}}
\def\Deformationcategory{\mathcal{D}\!\mathit{ef}}
\usepackage{cleveref}
\crefname{originallemma}{original lemma}{original lemmas}
\Crefname{originallemma}{Original lemma}{Original lemmas}
\makeatletter
\newlabel{derived-lemma-negative-vanishing}{{05TC}{source}{}{sourcepointer-1}{}}
\newlabel{derived-lemma-negative-vanishing@cref}{{[originallemma][0][]05TC}{[1][0][]source}{}{}{}}
\makeatother
\crefname{originallemma}{original lemma}{original lemmas}
\Crefname{originallemma}{Original lemma}{Original lemmas}
\makeatletter
\newlabel{derived-lemma-unbounded-right-derived}{{07K7}{source}{}{sourcepointer-2}{}}
\newlabel{derived-lemma-unbounded-right-derived@cref}{{[originallemma][0][]07K7}{[1][0][]source}{}{}{}}
\makeatother
\crefname{originallemma}{original lemma}{original lemmas}
\Crefname{originallemma}{Original lemma}{Original lemmas}
\makeatletter
\newlabel{etale-cohomology-lemma-cohomological-dimension-proper}{{095U}{source}{}{sourcepointer-3}{}}
\newlabel{etale-cohomology-lemma-cohomological-dimension-proper@cref}{{[originallemma][0][]095U}{[1][0][]source}{}{}{}}
\makeatother
\crefname{originallemma}{original lemma}{original lemmas}
\Crefname{originallemma}{Original lemma}{Original lemmas}
\makeatletter
\newlabel{etale-cohomology-lemma-proper-base-change-stalk}{{0DDF}{source}{}{sourcepointer-4}{}}
\newlabel{etale-cohomology-lemma-proper-base-change-stalk@cref}{{[originallemma][0][]0DDF}{[1][0][]source}{}{}{}}
\makeatother
\crefname{originallemma}{original lemma}{original lemmas}
\Crefname{originallemma}{Original lemma}{Original lemmas}
\makeatletter
\newlabel{etale-cohomology-lemma-relative-colimit-general}{{0EYM}{source}{}{sourcepointer-5}{}}
\newlabel{etale-cohomology-lemma-relative-colimit-general@cref}{{[originallemma][0][]0EYM}{[1][0][]source}{}{}{}}
\makeatother
\crefname{originaltheorem}{original theorem}{original theorems}
\Crefname{originaltheorem}{Original theorem}{Original theorems}
\makeatletter
\newlabel{etale-cohomology-theorem-colimit}{{09YQ}{source}{}{sourcepointer-6}{}}
\newlabel{etale-cohomology-theorem-colimit@cref}{{[originaltheorem][0][]09YQ}{[1][0][]source}{}{}{}}
\makeatother
\crefname{originallemma}{original lemma}{original lemmas}
\Crefname{originallemma}{Original lemma}{Original lemmas}
\makeatletter
\newlabel{lemma-base-change-f-shriek-separated}{{0F55}{source}{}{sourcepointer-7}{}}
\newlabel{lemma-base-change-f-shriek-separated@cref}{{[originallemma][0][]0F55}{[1][0][]source}{}{}{}}
\makeatother
\crefname{originallemma}{original lemma}{original lemmas}
\Crefname{originallemma}{Original lemma}{Original lemmas}
\makeatletter
\newlabel{lemma-colim-f-shriek-separated}{{0F54}{source}{}{sourcepointer-8}{}}
\newlabel{lemma-colim-f-shriek-separated@cref}{{[originallemma][0][]0F54}{[1][0][]source}{}{}{}}
\makeatother
\crefname{originallemma}{original lemma}{original lemmas}
\Crefname{originallemma}{Original lemma}{Original lemmas}
\makeatletter
\newlabel{lemma-f-shriek-composition}{{0F50}{source}{}{sourcepointer-9}{}}
\newlabel{lemma-f-shriek-composition@cref}{{[originallemma][0][]0F50}{[1][0][]source}{}{}{}}
\makeatother
\crefname{originallemma}{original lemma}{original lemmas}
\Crefname{originallemma}{Original lemma}{Original lemmas}
\makeatletter
\newlabel{lemma-lqf-f-shriek-separated-colimits}{{0F57}{source}{}{sourcepointer-10}{}}
\newlabel{lemma-lqf-f-shriek-separated-colimits@cref}{{[originallemma][0][]0F57}{[1][0][]source}{}{}{}}
\makeatother
\crefname{originallemma}{original lemma}{original lemmas}
\Crefname{originallemma}{Original lemma}{Original lemmas}
\makeatletter
\newlabel{lemma-proper-f-shriek}{{0F51}{source}{}{sourcepointer-11}{}}
\newlabel{lemma-proper-f-shriek@cref}{{[originallemma][0][]0F51}{[1][0][]source}{}{}{}}
\makeatother
\crefname{originallemma}{original lemma}{original lemmas}
\Crefname{originallemma}{Original lemma}{Original lemmas}
\makeatletter
\newlabel{morphisms-lemma-morphism-finite-type-bounded-dimension}{{0A3V}{source}{}{sourcepointer-12}{}}
\newlabel{morphisms-lemma-morphism-finite-type-bounded-dimension@cref}{{[originallemma][0][]0A3V}{[1][0][]source}{}{}{}}
\makeatother
\crefname{originalsection}{original section}{original sections}
\Crefname{originalsection}{Original section}{Original sections}
\makeatletter
\newlabel{section-compact-support}{{0F4W}{source}{}{sourcepointer-13}{}}
\newlabel{section-compact-support@cref}{{[originalsection][0][]0F4W}{[1][0][]source}{}{}{}}
\makeatother
\title{Base change map g\_! Rf'\_* K -> Rf\_*(g'\_! K) for proper f and its isomorphism in the torsion and bounded-below cases}
\author{The Stacks Project Authors}
\date{stacks-project:more-etale.tex:local-snapshot-20261012; GNU FDL 1.2 or later}
\begin{document}
\maketitle
\setcounter{subsection}{1}
\noindent\textit{Editorial note.} All delivered mathematical text is the author's own, copied from the pinned Stacks Project chapter file more-etale.tex; no statement or proof step was rewritten, reordered or removed. The counted span is the author environment and its own proof environment. References inside the span are handed to the assembler as navigation stubs that carry the canonical tag of the cited tag only; whether each of them is load-bearing is recorded on the stub itself.\par\medskip
\noindent\textit{Editorial note.} Disclosure: in case (a), the five-lemma verification between the displayed long exact sequences ends with 'Some details omitted.' (proof line 2815). The induction scheme, the distinguished triangle, and the map of triangles are fully displayed in the counted span, and case (b) is proved in full; the omitted step is the mechanical five-lemma diagram check on the displayed data.\par\medskip
\noindent\textit{Allowed standard prerequisites (editorial).}\begin{itemize}
\item Proper base change for stalks (Etale Cohomology, Lemma etale-cohomology-lemma-proper-base-change-stalk, canonical tag 0DDF). cited at proof lines 2738-2739 and 2781-2782. Identifies the stalks of R\textasciicircum{}q f\_* and kills the higher direct images in the torsion case.
\item Vanishing below the cohomological dimension (Derived Categories, Lemma derived-lemma-unbounded-right-derived, canonical tag 07K7). cited at proof lines 2729-2732. Reduces the case-(b) isomorphism to the vanishing of R\textasciicircum{}q f\_* on injective sheaves.
\end{itemize}\medskip
\ifdefined\hypertarget\hypertarget{sourcepointer-1}{}\fi
\noindent\textit{Source reference.} Original lemma 05TC: sources/stacks-chow/tags.txt, source lines 7476--7476. Cited at proof line 2794: both sides vanish for p below the vanishing bound in the induction start.\par\medskip
\ifdefined\hypertarget\hypertarget{sourcepointer-2}{}\fi
\noindent\textit{Source reference.} Original lemma 07K7: sources/stacks-chow/tags.txt, source lines 9647--9647. Cited at proof line 2729: reduces case (b) to vanishing of R\textasciicircum{}q f\_* on injective sheaves.\par\medskip
\ifdefined\hypertarget\hypertarget{sourcepointer-3}{}\fi
\noindent\textit{Source reference.} Original lemma 095U: sources/stacks-chow/tags.txt, source lines 11597--11597. Cited at proof lines 2727-2728: Rf\_* has finite cohomological dimension for proper f with bounded fibre dimension.\par\medskip
\ifdefined\hypertarget\hypertarget{sourcepointer-4}{}\fi
\noindent\textit{Source reference.} Original lemma 0DDF: sources/stacks-chow/tags.txt, source lines 16765--16765. Cited at proof lines 2738-2739 and 2781-2782: stalk identification of R\textasciicircum{}q f\_* and vanishing of higher direct images in the torsion case.\par\medskip
\ifdefined\hypertarget\hypertarget{sourcepointer-5}{}\fi
\noindent\textit{Source reference.} Original lemma 0EYM: sources/stacks-chow/tags.txt, source lines 18697--18697. Cited at proof line 2829: R\textasciicircum{}p f\_* commutes with the filtered colimit in the n-torsion filtration.\par\medskip
\ifdefined\hypertarget\hypertarget{sourcepointer-6}{}\fi
\noindent\textit{Source reference.} Original theorem 09YQ: sources/stacks-chow/tags.txt, source lines 12573--12573. Cited at proof lines 2764-2765: etale cohomology over X\_bar-y commutes with direct sums.\par\medskip
\ifdefined\hypertarget\hypertarget{sourcepointer-7}{}\fi
\noindent\textit{Source reference.} Original lemma 0F55: sources/stacks-chow/tags.txt, source lines 18925--18925. Cited at proof line 2741: formation of g'\_! commutes with base change.\par\medskip
\ifdefined\hypertarget\hypertarget{sourcepointer-8}{}\fi
\noindent\textit{Source reference.} Original lemma 0F54: sources/stacks-chow/tags.txt, source lines 18924--18924. Cited at proof line 2760: cohomology over the disjoint union of fibres splits as the direct sum.\par\medskip
\ifdefined\hypertarget\hypertarget{sourcepointer-9}{}\fi
\noindent\textit{Source reference.} Original lemma 0F50: sources/stacks-chow/tags.txt, source lines 18920--18920. Cited at proof lines 2701-2702: the second equality g\_! f'\_! = f\_! g'\_! in the map construction.\par\medskip
\ifdefined\hypertarget\hypertarget{sourcepointer-10}{}\fi
\noindent\textit{Source reference.} Original lemma 0F57: sources/stacks-chow/tags.txt, source lines 18927--18927. Cited at proof line 2827: g\_! and g'\_! commute with filtered colimits in the torsion-sheaf finish.\par\medskip
\ifdefined\hypertarget\hypertarget{sourcepointer-11}{}\fi
\noindent\textit{Source reference.} Original lemma 0F51: sources/stacks-chow/tags.txt, source lines 18921--18921. Cited at proof line 2700: the first and third equalities in the map construction.\par\medskip
\ifdefined\hypertarget\hypertarget{sourcepointer-12}{}\fi
\noindent\textit{Source reference.} Original lemma 0A3V: sources/stacks-chow/tags.txt, source lines 12753--12753. Cited at proof lines 2723-2725: localize on Y so that the fibre dimension of f is bounded.\par\medskip
\ifdefined\hypertarget\hypertarget{sourcepointer-13}{}\fi
\noindent\textit{Source reference.} Original section 0F4W: sources/stacks-chow/tags.txt, source lines 18916--18916. Cited at proof line 2762: the definition of g'\_\{bar-y,!\} used in the disjoint-fibre step.\par\medskip
\begin{lemma}
\label{lemma-shriek-proper-and-open}
Consider a commutative diagram of schemes
$$
\xymatrix{
X' \ar[r]_{g'} \ar[d]_{f'} & X \ar[d]^f \\
Y' \ar[r]^g & Y
}
$$
with $f$ and $f'$ proper and $g$ and $g'$ separated and locally quasi-finite.
Let $\Lambda$ be a ring. Functorially in $K \in D(X'_\etale, \Lambda)$
there is a canonical map
$$
g_!Rf'_*K \longrightarrow Rf_*(g'_!K)
$$
in $D(Y_\etale, \Lambda)$. This map is an isomorphism if
(a) $K$ is bounded below and has torsion cohomology sheaves, or
(b) $\Lambda$ is a torsion ring.
\end{lemma}

\begin{proof}
Represent $K$ by a K-injective complex $\mathcal{J}^\bullet$ of
sheaves of $\Lambda$-modules on $X'_\etale$. Choose a quasi-isomorphism
$g'_!\mathcal{J}^\bullet \to \mathcal{I}^\bullet$ to a
K-injective complex $\mathcal{I}^\bullet$ of
sheaves of $\Lambda$-modules on $X_\etale$.
Then we can consider the map
$$
g_!f'_*\mathcal{J}^\bullet =
g_!f'_!\mathcal{J}^\bullet =
f_!g'_!\mathcal{J}^\bullet =
f_*g'_!\mathcal{J}^\bullet \to
f_*\mathcal{I}^\bullet
$$
where the first and third equality come from
Lemma \ref{lemma-proper-f-shriek}
and the second equality comes from
Lemma \ref{lemma-f-shriek-composition} which tells us that both
$g_! \circ f'_!$ and $f_! \circ g'_!$ are equal to
$(g \circ f')_! = (f \circ g')_!$ as subsheaves of
$(g \circ f')_* = (f \circ g')_*$.

\medskip\noindent
Assume $\Lambda$ is torsion, i.e., we are in case (b).
With notation as above, it suffices to show that
$f_*g'_!\mathcal{J}^\bullet \to f_*\mathcal{I}^\bullet$
is an isomorphism. The question is local on $Y$.
Hence we may assume that the dimension
of fibres of $f$ is bounded, see Morphisms, Lemma
\ref{morphisms-lemma-morphism-finite-type-bounded-dimension}.
Then we see that $Rf_*$ has finite cohomological dimension, see
\'Etale Cohomology, Lemma
\ref{etale-cohomology-lemma-cohomological-dimension-proper}.
Hence by Derived Categories, Lemma \ref{derived-lemma-unbounded-right-derived},
if we show that $R^qf_*(g'_!\mathcal{J}) = 0$ for $q > 0$
and any injective sheaf of $\Lambda$-modules $\mathcal{J}$
on $X'_\etale$, then the result follows.

\medskip\noindent
The stalk of $R^qf_*(g'_!\mathcal{J})$ at a geometric point
$\overline{y}$ is equal to
$H^q(X_{\overline{y}}, (g'_!\mathcal{J})|_{X_{\overline{y}}})$ by
\'Etale Cohomology, Lemma
\ref{etale-cohomology-lemma-proper-base-change-stalk}.
Since formation of $g'_!$ commutes with base change
(Lemma \ref{lemma-base-change-f-shriek-separated})
this is equal to
$$
H^q(X_{\overline{y}}, g'_{\overline{y}, !}(\mathcal{J}|_{X'_{\overline{y}}}))
$$
where $g'_{\overline{y}} : X'_{\overline{y}} \to X_{\overline{y}}$
is the induced morphism between geometric fibres.
Since $Y' \to Y$ is locally quasi-finite, we see that
$X'_{\overline{y}}$ is a disjoint union of the fibres
$X'_{\overline{y}'}$ at geometric points $\overline{y}'$ of $Y'$
lying over $\overline{y}$. Denote
$g'_{\overline{y}'} : X'_{\overline{y}'} \to X_{\overline{y}}$
the restriction of $g'_{\overline{y}}$ to $X'_{\overline{y}'}$.
Thus the previous cohomology group is equal to
$$
H^q(X_{\overline{y}},
\bigoplus\nolimits_{\overline{y}'/\overline{y}}
g'_{\overline{y}', !}(\mathcal{J}|_{X'_{\overline{y}'}}))
$$
for example by Lemma \ref{lemma-colim-f-shriek-separated} (but it
is also obvious from the definition of $g'_{\overline{y}, !}$
in Section \ref{section-compact-support}).
Since taking \'etale cohomology over $X_{\overline{y}}$
commutes with direct sums
(\'Etale Cohomology, Theorem \ref{etale-cohomology-theorem-colimit})
we conclude it suffices to show that
$$
H^q(X_{\overline{y}}, g'_{\overline{y}', !}(\mathcal{J}|_{X'_{\overline{y}'}}))
$$
is zero. Observe that
$g_{\overline{y}'} : X'_{\overline{y}'} \to X_{\overline{y}}$
is a morphism between proper scheme over $\overline{y}$ and hence is
proper itself. As it is locally quasi-finite as well we conclude that
$g_{\overline{y}'}$ is finite. Thus we see that
$g'_{\overline{y}', !} = g'_{\overline{y}', *} = Rg'_{\overline{y}', *}$.
By Leray we conlude that we have to show
$$
H^q(X'_{\overline{y}'}, \mathcal{J}|_{X'_{\overline{y}'}})
$$
is zero. As $\Lambda$ is torsion, this follows from proper base change
(\'Etale Cohomology, Lemma
\ref{etale-cohomology-lemma-proper-base-change-stalk})
as the higher direct images of $\mathcal{J}$ under $f'$ are zero.

\medskip\noindent
Proof in case (a). We will deduce this from case (b) by standard arguments.
We will show that the induced map $g_! R^pf'_* K \to R^pf_*(g'_!K)$
is an isomorphism for all $p \in \mathbf{Z}$. Fix an integer
$p_0 \in \mathbf{Z}$. Let $a$ be an integer such that $H^j(K) = 0$
for $j < a$. We will prove $g_! R^pf'_* K \to R^pf_*(g'_!K)$
is an isomorphism for $p \leq p_0$ by descending induction on $a$. If
$a > p_0$, then we see that the left and right hand side of
the map are zero for $p \leq p_0$ by trivial vanishing, see
Derived Categories, Lemma \ref{derived-lemma-negative-vanishing}
(and use that $g_!$ and $g'_!$ are exact functors).
Assume $a \leq p_0$. Consider the distinguished triangle
$$
H^a(K)[-a] \to K \to \tau_{\geq a + 1}K
$$
By induction we have the result for $\tau_{\geq a + 1}K$.
In the next paragraph, we will prove the result for $H^a(K)[-a]$.
Then five lemma applied to the map between
long exact sequence of cohomology sheaves
associated to the map of distinguished triangles
$$
\xymatrix{
g_! Rf'_*(H^a(K)[-a]) \ar[d] \ar[r] &
g_! Rf'_* K \ar[r] \ar[d] &
g_! Rf'_* \tau_{\geq a + 1} K \ar[d] \\
Rf_*(g'_!(H^a(K)[-a])) \ar[r] &
Rf_*(g'_!K) \ar[r] &
Rf_*(g'_!\tau_{|geq a + 1}K)
}
$$
gives the result for $K$. Some details omitted.

\medskip\noindent
Let $\mathcal{F}$ be a torsion abelian sheaf on $X'_\etale$.
To finish the proof we show that
$g_! Rf'_*\mathcal{F} \to R^pf_*(g'_!\mathcal{F})$
is an isomorphism for all $p$.
We can write $\mathcal{F} = \bigcup \mathcal{F}[n]$ where
$\mathcal{F}[n] = \Ker(n : \mathcal{F} \to \mathcal{F})$.
We have the isomorphism for $\mathcal{F}[n]$ by case (b).
Since the functors $g_!$, $g'_!$, $R^pf_*$, $R^pf'_*$ commute
with filtered colimits (follows from
Lemma \ref{lemma-lqf-f-shriek-separated-colimits} and
\'Etale Cohomology, Lemma
\ref{etale-cohomology-lemma-relative-colimit-general})
the proof is complete.
\end{proof}

\end{document}
